% Supplementary Material, rater pool, agreement strata and power (compiled only into the arXiv edition via 00_main_arxiv.tex).
% Provenance: experiments/44_* (RESULTS.md, results/key_numbers.md, results/*.csv), experiments/47_* (results/label_budget_summary.json), experiments/48_* (README.md tables), experiments/50_* (README.md, results/strata_tables.md), experiments/52_* (README.md); Lean theorem names from docs/notes/point2_support_27Nov2026/lean (README.md claim map, axioms_output_27Nov2026.txt).
% One sentence per line (Trigger Rule 13); strictly anonymized.
% Notation follows the main text: $x$ is the input, $r$ the human rating, $p^*(r\mid x,\mathcal{P},\rho)$ the human rating distribution for a stated rater pool $\mathcal{P}$ and rubric $\rho$, $q(r\mid x)$ the judge's forecast.

\section{Rater Pool, Agreement Strata and Power}
\label{sec:supp-pool-strata}

This section reports the analyses that bear on the target $p^*(r\mid x,\mathcal{P},\rho)$ and on the stratified protocol of the main text.
Except for the linear heads of Experiment~50 and the simulation of Experiment~52, they score five frozen frontier judges on public multi-annotator rating sets with nothing fitted; the API cost is zero, and the Experiment~44 protocol was committed before any analysis ran.
Scores are the multiclass Brier score (ChaosNLI; unnormalized, twice the main text's normalized score) and the ranked probability score (RPS) normalized by $K-1$ (DICES), both lower is better, computed as the per-sample risk defined in the main text.
As a reference, a uniform forecast scores $0.6667$ on three-class Brier and $0.2778$ or $0.1111$ on three-class normalized RPS, depending on the rating.

\paragraph{Does the rater pool change the verdict?}
In the two DICES sets we replaced the pool that defines $p^*$ by demographic sub-pools (gender, age group and race), with $11$ to $139$ raters per pool.
The item-level rating distributions of two pools differ by a mean total variation distance of $0.08$ to $0.18$, but so do two random rater groups, and no contrast is distinguishable from a rater-label permutation null (smallest Holm-adjusted $p=0.59$ over six contrasts).
No judge-by-contrast score difference is Holm-significant ($0$ of $30$ cells), and the Kendall $\tau_b$ between the judges' orders under the two pools is $0.8$ or $1.0$ in all six contrasts; the best and second-best judge keep their places in all twelve orders.
We read this as a statement about power, not as evidence that the pool never matters: the smallest pools have $11$ and $27$ raters, and the judges were asked to forecast a whole crowd, not a sub-pool.
The construct statement of the main text is therefore a requirement of the method, not a finding of this analysis.

\paragraph{Consensus and contested items.}
Agreement is measured on one random half of each item's ratings and the other half is scored; the strata are the lowest and the highest third of items by agreement.
Sharpening a judge's forecast (raising it to the power $1/\tau$ with $\tau=0.5$) helps relatively more, or hurts less, on consensus items than on contested items in $25$ of $25$ judge-by-set cells, and the item-bootstrap interval of that difference (strata resampled independently) excludes $0$ in all $25$.
A peaked forecast is therefore penalized much less where raters agree than where they disagree, which is why scores are reported by stratum.
The picture is narrower than ``judges should be peaked on consensus items'': sharpening improved the consensus stratum in only $5$ of $25$ cells, and on contested items the softest temperature on the grid ($\tau=2$) was best in $23$ of $25$, so these judges are overconfident there.
Stratifying five ways by agreement (Experiment~50), the skill of the frontier judges over the base rate falls from $0.27$--$0.91$ in Brier skill on the three ChaosNLI sets and $0.16$--$0.30$ in normalized-RPS skill on the two DICES sets in the highest-agreement fifth of items to $-0.28$--$0.11$ and $-0.38$--$0.04$ in the lowest, while their Pearson correlation on the contested NLI items stays at $0.38$--$0.86$.
A pooled correlation hides where a judge stops being useful; a skill score reported by agreement stratum shows it.
In the same experiment the penalty for training a linear head on correlation instead of cross-entropy was significantly smaller in the lowest-agreement bin than in the highest on $9$ of $12$ views (a $95\%$ item-bootstrap interval below zero in at least three of five rating splits, with no correction for multiplicity; three of the nine are SummEval expert views, binned and scored on the same ratings), contrary to our frozen prediction, so we do not claim that correlation-trained judges fail mainly on ambiguous items.

\paragraph{Binned divergence.}
Binned Jensen--Shannon divergence~\citep[their Eqs.~1 and~2]{elangovan2025beyond}, adapted to probabilistic forecasts by replacing their pooled sampled judge labels with the mean forecast, compares the pooled human ratings with that mean forecast inside bins of items defined by the human median or the most frequent human label.
It is exactly unchanged when each judge's forecasts are reassigned among the items of the same bin (largest change $2.4\times10^{-16}$ over $200$ shuffles in each of $25$ judge-by-set cells), while the proper-score risk worsens in all $25$ cells, by $0.0571$ to $0.1075$ in Brier score on the three ChaosNLI sets and $0.0086$ to $0.0202$ in normalized RPS on the two DICES sets.
Binned divergence cannot see which forecast belongs to which item inside a bin, which is the shortcoming its authors name.
The top judge differs between binned divergence and the proper score in $2$ of $5$ sets; with five judges and no test, we report that as descriptive only.

\paragraph{How many items does a paired comparison need at one rating per item?}
With $N$ held-out items and $k$ ratings each, the variance of a paired mean difference is $V(N,k)=(\sigma_b^2+\sigma_w^2/k)/N$ for between-item and within-item variance components estimated from the data.
At $N=400$ and $k=1$ the minimum detectable paired difference at $80\%$ power is $0.034$ to $0.065$ in Brier score on the three ChaosNLI sets and $0.014$ to $0.025$ in normalized RPS on the two DICES sets, and $0$ to $6$ of the $10$ judge pairs per set are resolvable.
At a fixed total number of ratings the variance never decreases in $k$, so more items beat more ratings per item, and a synthetic check of the skill estimator (Experiment~52) agrees: at $500$ ratings the standard deviation is $0.029$ with one rating on each of many items against $0.035$ with five ratings on a fifth as many.
Nearly tied judges need far more items than these sets contain ($N_{80}$ up to $1.9\times10^{6}$ for the nearest ties).

\paragraph{Calibration needs in-domain held-out ratings.}
For the frontier judges (Experiment~56), calibrating each judge on $n$ held-out items with one rating each left its mean skill below its uncalibrated skill in $25$ of $25$ judge-set cases at $n=10$, $21$ of $25$ at $n=25$, $12$ of $25$ at $n=50$ and $8$ of $25$ at $n=100$; $n=100$ is the smallest tested budget at which every judge's mean calibrated skill is above the base rate on all five sets, and at $n=200$ it was still $0.003$ to $0.022$ below a map fitted on all ratings of the pool.
For the linear heads of Experiment~41, a circular comparison kept here as a check only (both heads are scored by a proper rule), the largest five-seed mean gap of $0.0024$ normalized RPS between the cross-entropy-trained and the Pearson-trained head is measured after each head's score passes through its own map, fitted by cross-entropy on the full validation split of the same data.
Fitting that map on $n=25$ single held-out ratings instead (Experiment~47), the calibrated Pearson-trained head is still worse than the raw cross-entropy-trained head on all sixteen views, by $0.0097$ to $0.0225$ normalized RPS, and it reaches the $0.0024$ bound only at $n=50$ (one view), $n=100$ (six views) or $n=200$ (nine views).
Fitting the map on a different dataset or rubric (Experiment~48) reopens the proper-score gap in $53$ of $82$ (judge, source, target) cases under the frozen rule (five frozen frontier judges and one G-Eval judge; $33$ of the $53$ move a map between NLI and DICES, which differ in label set and rubric, so those cases also reflect a label mismatch); the transferred map is significantly worse than the raw forecast in $33$ of the $82$ cases ($21$ of them between NLI and DICES) and significantly worse than a map fitted in domain in $78$ of the $82$.
For the linear heads the gap reopens in all $16$ leave-one-dataset-out views and in $5$ of $6$ expert-to-crowd rater-pool shifts; the shift between the two DICES sets is the exception ($0$ of $10$ cases reopen).
These results concern the cost of calibration, not an argument against it: each new target needs its own held-out human ratings, scored with a proper rule, which is the recommendation of this paper.

\paragraph{Machine-checked claims.}
The identification claims are proved in Lean~4 for finite laws (Lean 4.36.0, Mathlib revision \texttt{fc4f59e}, zero \texttt{sorry}, only the standard axioms \texttt{propext}, \texttt{Classical.choice} and \texttt{Quot.sound}).
Forecasts with the same conditional mean tie on Pearson, Spearman, Kendall and scalar mean squared error of the mean readout (\texttt{Point2.FinLaw.same\_\allowbreak mean\_\allowbreak tie}, \texttt{Point2.scalar\_\allowbreak same\_\allowbreak mean\_\allowbreak tie}).
The Brier score and RPS are strictly proper (\texttt{Point2.brier\_\allowbreak strictly\_\allowbreak proper}, \texttt{Point2.rps\_\allowbreak strictly\_\allowbreak proper}); a scalar encoding is strictly proper for two categories and not for three or more (\texttt{Point2.scalar\_\allowbreak two\_\allowbreak strictly\_\allowbreak proper}, \texttt{Point2.scalar\_\allowbreak not\_\allowbreak strictly\_\allowbreak proper}); and squared correlation is one minus the mean squared error of the best affine fit divided by the variance of the rating (\texttt{Point2.FinLaw.r\_\allowbreak sq\_\allowbreak identity}).
For a finite joint law, a collapsed judge that reports an increasing affine function of the true conditional mean has correlation with the human rating equal to the square root of the correlation of a truthful judge that samples from the true conditional distribution independently of the human rating, and the square root is strictly larger when that correlation lies strictly between $0$ and $1$ (\texttt{Point2.Sampled.Model.anti\_\allowbreak proper}).
Cross-entropy strict propriety is cited (Gibbs' inequality), not formalized.
That agreement with a label depends on a forecast with a strict mode only through that mode is proved in Lean~4 without Mathlib for integer-scaled (that is, rational) forecasts and any number of categories (\texttt{Point1.agrees\_\allowbreak iff\_\allowbreak mode}, \texttt{Point1.agreement\_\allowbreak depends\_\allowbreak only\_\allowbreak on\_\allowbreak mode}).
Five closed three-category instances are compiled as well: two forecasts with the same mode tie on agreement while the Brier excess separates them, two forecasts with the same mean differ in agreement, a change of rater pool changes the Brier excess of a fixed forecast, and sharpening is rewarded on a point-mass item and penalized on a contested one (\texttt{Point1.same\_\allowbreak mode\_\allowbreak ties\_\allowbreak agreement\_\allowbreak but\_\allowbreak brier\_\allowbreak separates\_\allowbreak instance}, \texttt{Point1.same\_\allowbreak mean\_\allowbreak different\_\allowbreak agreement\_\allowbreak instance}, \texttt{Point1.pool\_\allowbreak changes\_\allowbreak the\_\allowbreak ruler\_\allowbreak instance}, \texttt{Point1.consensus\_\allowbreak item\_\allowbreak rewards\_\allowbreak peaked\_\allowbreak forecast\_\allowbreak instance}, \texttt{Point1.contested\_\allowbreak item\_\allowbreak penalizes\_\allowbreak sharpening\_\allowbreak instance}).
A seeded exact-fraction script checks the same claims on random finite instances, including the claims that must fail: forecasts with the same mode need not have the same Brier excess, and forecasts with the same mean need not have the same agreement.
Kappa statements are paper proofs, not yet Lean-checked; the ratio bias of the skill score is checked numerically (Experiment~52), not in Lean.
